\lab{17}{Автотесты (юнит-тесты)}
\task
Подготовить воспроизводимые модульные тесты основных компонент
поисковой системы. Проверить штатные и граничные случаи,
некорректные входы, согласованность форматов и результатов.

\method
Используется Google Test, вызывающий функции C++ напрямую.
CMake подключает установленный пакет или закреплённую версию
1.15.2 с проверкой SHA-256. CTest регистрирует отдельные тесты,
поэтому отчёт показывает конкретную проверку, вызвавшую ошибку.
Тесты запускаются отдельно от обычной сборки.

Искусственный корпус содержит восемь документов:
\begin{center}
\begin{tabular}{rl}\toprule
ID & Текст\\\midrule
0 & \code{alpha beta beta gamma}\\
1 & \code{alpha x beta gamma}\\
2 & \code{beta alpha gamma}\\
3 & \code{alpha beta gamma}\\
4 & пустой документ\\
5 & \code{delta gamma}\\
6 & \code{alpha alpha beta}\\
7 & \code{common common delta}\\\bottomrule
\end{tabular}
\end{center}
Ручной эталон фиксирует ответ независимо от поискового кода.
Например, \code{"alpha beta"} даёт \(\{0,3,6\}\),
\code{"alpha beta" / 2} - \(\{0,1,3,6\}\),
\code{!alpha} - \(\{4,5,7\}\).
Один набор запросов выполняется по raw, VB и индексу,
построенному с очень малым буфером, с прыжками и без них.

\journal
Сначала добавлены тесты кодека и булевых выражений. Затем они
расширены цитатами, внешним слиянием, метаданными и нормами TF-IDF.
Для каждого случая используется временный каталог, а данные
создаются самим тестом: рабочий корпус и готовые индексы
не перезаписываются. \code{SCOPED_TRACE} добавляет к диагностике
текущий запрос. Для псевдослучайных целых кодека задано фиксированное зерно.

\testing
\begin{longtable}{p{0.41\textwidth}p{0.51\textwidth}}
\toprule Проверка & Что обнаруживает\\\midrule
\code{BoundariesAndInvalidEncoding} & Ошибки границ LEB128, каноничности и восстановления целых.\\
\code{BooleanAndOrderedPhrases} & Ошибки приоритетов, цитат, порядка, отрицания и повторных слов.\\
\code{MergeMetadataNormsAndSkips} & Зависимость файла от размера блока, повреждение метаданных, норм и курсоров.\\
\code{FuzzySearchAndTfIdfRanking} & Нарушение строгих фильтров, частичных совпадений и ранжирования.\\\bottomrule
\end{longtable}
Ручная норма первого документа использует
\(idf=1+\ln(9/6)\), \(\beta=1+\ln2\),
\(\|d_0\|=idf\sqrt{2+\beta^2}\).
Пустой документ имеет нулевые длину и норму. Побайтовое
сравнение VB-файлов при разном размере буфера проверяет,
что внешнее слияние не меняет результат сборки.

\results
Повторно запущены модульные тесты и проверки командного интерфейса.
Тесты создают только небольшие временные данные, полные индексы
не перестраиваются. Получены результаты:
\begin{table}[H]\centering
\begin{tabular}{lrr}\toprule
Проект & Пройдено & Ошибок\\\midrule
Оценка качества, №6 & 3 & 0\\
Булев индекс, №7 & 2 & 0\\
Булев поиск, №8 & 2 & 0\\
Координатный поиск, №10 & 4 & 0\\
Сжатие, №11 & 4 & 0\\
Прыжки, №12 & 4 & 0\\
TF-IDF, №13 & 4 & 0\\
Общий набор, №17 & 4 & 0\\
Проверка оценочного конвейера & 7 & 0\\
Регрессионные проверки CLI & 6 & 0\\\midrule
Итого запущено & 40 & 0\\\bottomrule
\end{tabular}
\caption{Результаты C++- и Python-тестов. Наборы 10--13 и 17 частично дублируются.}
\end{table}
Фрагмент журнала общего набора:
\begin{lstlisting}
[ RUN      ] Codec.BoundariesAndInvalidEncoding
[       OK ] Codec.BoundariesAndInvalidEncoding
[ RUN      ] PositionalIndex.BooleanAndOrderedPhrases
[       OK ] PositionalIndex.BooleanAndOrderedPhrases
[ RUN      ] PositionalIndex.MergeMetadataNormsAndSkips
[       OK ] PositionalIndex.MergeMetadataNormsAndSkips
[ RUN      ] PositionalIndex.FuzzySearchAndTfIdfRanking
[       OK ] PositionalIndex.FuzzySearchAndTfIdfRanking
[  PASSED  ] 4 tests.
\end{lstlisting}
Дополнительные семь тестов оценивания проверяют валидность
и полноту разметки. Шесть регрессионных CLI-тестов проверяют
повторные измерения без корпуса, неизменность индексов,
ошибки при отсутствии подготовленных файлов, отказ от повторной
записи существующего индекса и точность преобразования шага прыжков.
Для сравнения индексов используются SHA-256 и побайтовая проверка.

Число 40 означает количество выполненных тестов, а не число
независимых алгоритмов: часть C++-наборов повторяется в разных проектах.
Автотесты робота, Python-обвязки токенизатора и нормализатора
в этот итог не включены.

\conclusions
Малый корпус с ручными ответами позволяет проверять семантику
и ловить регрессии при смене кодека и способа сборки. Тесты
не доказывают отсутствие ошибок и не являются бенчмарком.
HTTP-интерфейс и настоящая MongoDB требуют отдельных интеграционных
проверок. Длительные сбои диска и сети требуют нагрузочных и аварийных
испытаний. Релевантность оценивается на размеченном пуле, а не
модульными тестами алгоритма.
